\documentclass[12pt]{article}
\usepackage[utf8]{inputenc}
\usepackage[T1]{fontenc}
\usepackage{hyperref}
\title{Rinc\'on Zapotec Language Resources for K-12 Curriculum Development}
\author{Compiled by Claude Sonnet 5.5 (claude-sonnet-5-5) for Andr\'as Kornai}
\date{}
\begin{document}
\maketitle

\section{Introduction}
Rinc\'on Zapotec (ISO 639-3 zar; Glottocode rinc1236) is an Otomanguean language of Oaxaca, Mexico, with about 41,000 speakers according to the reference description. The searches run for this report found very little specific to this language, and it must not be confused with Southern Rinc\'on Zapotec (zsr) or Sierra de Ju\'arez Zapotec (zaa), which have separate codes.

\section{Primary Educational Resources}
No primers, textbooks or school programmes for Rinc\'on Zapotec were found.

\section{Linguistic and Reference Materials}
No dictionaries, grammars or academic descriptions of zar were found in the search results.

\section{Religious and Community Resources}
No Scripture or community resources specific to zar were confirmed. The neighbouring but distinct Sierra de Ju\'arez Zapotec (zaa) has Scripture and Jesus Film materials, but these are not for zar.

\section{Digital and Audio Resources}
Global Recordings Network has a catalogue page for Rinc\'on Zapotec: \url{https://globalrecordings.net/en/language/zar} The search snippet did not show what recordings it holds, so its content should be checked directly.

\section{Curriculum Development Context}
Little is documented online for this variety. Curriculum work would need to rely on speaker communities and on specialist literature searched in Spanish, and care should be taken to keep zar separate from zsr and zaa.

\end{document}